\chapter{Óptica Ondulatoria}
\label{ch:ondulatoria}


La visión geométrica de la luz, si bien elegante y suficiente para muchas de las aplicaciones elementales, no capta por completo la verdadera naturaleza de la luz, pues en la naturaleza la luz no se mantiene confinada en un rayo mientras se propaga por el espacio libre, sino que tiende a ``esparcirse'' hacia las vecindades, entre otros fenómenos que la teoría geométrica no explica ni contempla.

En concreto, esta manera geométrica de entender la luz se vio enfrentada a serios problemas para explicar fenómenos como la interferencia, la difracción y la birrefringencia, por mencionar algunos. Con el tiempo se fueron acumulando experiencias que contradecían el modelo geométrico, hasta que las falencias de este modelo se hicieron innegables.

Por otro lado, aparentemente inconexo en un principio, las investigaciones concernientes a los fenómenos eléctricos y magnéticos dieron fruto a grandes avances en el entendimiento de estas fuerzas a distancia observadas en cuerpos ``terrenales'' como la magnetita. Fue Faraday quien entendió cómo se influían mutuamente, lo que sugería una misma naturaleza, y también observó experimentalmente cómo la luz era influida por estas fuerzas, ahora mejor entendidas. James Clerk Maxwell, en una tarea verdaderamente unificadora y excelsa matemáticamente, sintetizó las ideas de su época en las ecuaciones que llevan su nombre y, por una necesidad física de conservación y el conocimiento matemático de las ecuaciones que la aseguraban, fue más allá que sus predecesores postulando las ondas como consecuencia de respetar este principio de conservación para la carga eléctrica \cite{maxwell1865}.

Las ondas electromagnéticas de Maxwell tenían que viajar a cierta velocidad que resultó ser la misma que la velocidad a la que viajaban las ondas luminosas. 
Finalmente, se llegó a concluir que la luz visible era radiación electromagnética. 



\section{Ecuaciones de Maxwell}

En su forma macroscópica, las ecuaciones de Maxwell en un medio material pueden escribirse como \cite{jackson1999}
\begin{align}
\nabla \cdot \mathbf{D}(\mathbf x, t) &= \rho_{\ell}(\mathbf x, t)
& \qquad
\nabla \cdot \mathbf{B}(\mathbf x, t) &= 0, \\
\nabla \times \mathbf{E}(\mathbf x, t) &= -\,\frac{\partial \mathbf B(\mathbf x, t)}{\partial t}
& \qquad
\nabla \times \mathbf H(\mathbf x, t) &= \mathbf J_{\ell}(\mathbf x, t)
+ \frac{\partial \mathbf D(\mathbf x, t)}{\partial t},
\label{eq:maxwell_equaitons}
\end{align}
donde $\rho_{\ell}$ y $\mathbf J_{\ell}$ son las densidades de carga y de corriente \emph{libres}, y los campos materiales $\mathbf D$ y $\mathbf H$ absorben en su definición las cargas y corrientes \emph{ligadas} del medio, asociadas a la polarización y a la magnetización \cite{jackson1999}.

El interés de este trabajo es la propagación en \emph{dieléctricos simples}: medios lineales, homogéneos e isótropos, para los cuales las relaciones constitutivas son
\[
\mathbf D = \epsilon\,\mathbf E, \qquad \mathbf B = \mu\,\mathbf H,
\]
con permitividad $\epsilon$ y permeabilidad $\mu$ constantes. En estos medios no existen cargas ni corrientes libres,
\[
\rho_{\ell}(\mathbf x,t) = 0, \qquad \mathbf J_{\ell}(\mathbf x,t) = \mathbf 0,
\]
aunque sí existen cargas y corrientes ligadas, cuya respuesta queda absorbida en los campos materiales a través de $\epsilon$ y $\mu$. Con esto, las ecuaciones toman la forma
\begin{align}
\nabla \cdot \mathbf{E}(\mathbf x, t) &= 0
& \qquad
\nabla \cdot \mathbf{B}(\mathbf x, t) &= 0,
\\
\nabla \times \mathbf{E}(\mathbf x, t) &= -\,\frac{\partial \mathbf B(\mathbf x, t)}{\partial t}
& \qquad
\nabla \times \mathbf B(\mathbf x, t) &= \mu\epsilon\,\frac{\partial \mathbf E(\mathbf x, t)}{\partial t}.
\label{eq:maxwell_equation_simple_dielectric}
\end{align}

De la combinación adecuada de estas ecuaciones se concluye que los campos satisfacen, en el dieléctrico simple, ecuaciones de onda de la forma
\begin{align}
\nabla^2 \mathbf E(\mathbf x, t) - \mu\epsilon\,\frac{\partial^2 \mathbf E(\mathbf x, t)}{\partial t^2} &= 0,
\label{eq:wave_equation_E}
\\
\nabla^2 \mathbf B(\mathbf x, t) - \mu\epsilon\,\frac{\partial^2 \mathbf B(\mathbf x, t)}{\partial t^2} &= 0,
\label{eq:wave_equation_B}
\end{align}
donde
\[
v = \frac{1}{\sqrt{\mu \epsilon}} = \frac{c}{n}
\]
es la velocidad de propagación de la onda electromagnética en el medio, con
$c = 1/\sqrt{\mu_0\epsilon_0}$ la velocidad de la luz en el vacío y
$n = \sqrt{\mu\epsilon/\mu_0\epsilon_0}$ el índice de refracción del dieléctrico \cite{jackson1999, bornwolf1999}.

Por el teorema de Fourier sabemos que toda función admite una representación armónica \cite{arfken2013}, de modo que podemos escribir el campo eléctrico en dicha representación como
\begin{equation}
    \vb{E}(\vb x, t) = \frac{1}{2\pi}\int_{-\infty}^{\infty} \vb{E}_\omega(\vb x) e^{-i\omega t} d\omega.
    \label{eq:E_field_harmonic_representation}
\end{equation}
En la representación armónica el campo magnético se obtiene fácilmente, como
\begin{equation}
    \vb{B}_\omega = \frac{1}{i\omega} \nabla \times \vb{E}_\omega,
    \label{eq:B_harmonic_func_of_E_harmonic}
\end{equation}
por lo que en lo sucesivo omitiremos el campo magnético en nuestro análisis, pues la información del campo eléctrico nos es suficiente. De este modo, cada componente armónica satisface
\begin{equation}
    \nabla^2 \vb{E}(\vb x) + k^2 \vb E = 0,
    \label{eq:E_field_wave_equation}
\end{equation}
donde
\[
k = \frac{\omega}{v} = \omega\sqrt{\mu\epsilon} = \frac{n\omega}{c}
\]
es el número de onda en el dieléctrico; esta es la ecuación fundamental para la propagación del campo armónico en el medio.

La ecuación \eqref{eq:E_field_wave_equation} es equivalente a una ecuación de Helmholtz escalar para cada una de sus componentes en una base independiente del espacio; sea $\Psi(\vb x)$ una de esas componentes, escribimos
\begin{equation}
    \nabla^2 \Psi (\vb x) + k^2 \Psi (\vb x) = 0,
    \label{eq:Helmholtz_equation}
\end{equation}
que es la ecuación base de los \textit{métodos escalares de propagación}. Estos métodos son, en últimas, soluciones a la ecuación escalar de Helmholtz dadas unas condiciones de contorno, y se dividen en exactos y aproximados según si resuelven la ecuación \eqref{eq:Helmholtz_equation} de forma exacta o aproximada \cite{goodman2005}.


% ======================
\section{Propagación de los campos}

En la sección anterior partimos de la ecuación vectorial de Helmholtz,
ecuación~\eqref{eq:E_field_wave_equation}, y observamos que, al expresar el
campo en una base lineal ---entendida aquí como una base independiente de la
posición---, cada componente satisface por separado la ecuación escalar de Helmholtz,
ecuación~\eqref{eq:Helmholtz_equation}. Esto permite emplear los métodos
escalares de difracción para propagar las componentes del campo dentro de un
medio homogéneo.

Sin embargo, satisfacer la ecuación de Helmholtz componente a componente es
una condición necesaria, pero no suficiente para obtener un campo
electromagnético admisible: el campo eléctrico también debe cumplir la
condición de transversalidad
\begin{equation}
\nabla\cdot\vb E=0.
\label{eq:transversality_propagation}
\end{equation}
Tomando \(z\) como la dirección longitudinal, escribimos
\begin{equation}
\vb E
=
\vb E_\perp+E_z\uvec z.
\end{equation}
El método consiste en propagar escalarmente las dos componentes
transversales,
\begin{equation}
\vb E_\perp(\vb r,z)
=
\sum_{j=1}^{2}
\mathscr R(z)\!\left[E_j(\vb r,0)\right]\uvec{e}_j,
\end{equation}
donde \(\mathscr R(z)\) es el operador que, dado el campo en el plano inicial,
determina el campo en un plano paralelo situado a una distancia \(z\), de
acuerdo con la ecuación de Helmholtz. La condición de transversalidad
determina entonces la componente longitudinal mediante
\begin{equation}
E_z(\vb r,z)
=
-\mathscr F_\perp^{-1}
\left\{
\frac{\vb k_\perp}{k_z}\cdot
\mathscr F_\perp\!\left\{\vb E_\perp(\vb r,z)\right\}
\right\},
\qquad
k_z=\sqrt{k^2-\lVert\vb k_\perp\rVert^2},
\label{eq:Ez_from_transverse_spectrum}
\end{equation}
donde \(\vb k_\perp\) es la componente transversal del vector de onda, y
\(\mathscr F_\perp\) y \(\mathscr F_\perp^{-1}\) denotan las transformadas de
Fourier transversal directa e inversa, respectivamente. Esta reconstrucción se
aplicará explícitamente en el capítulo~\ref{ch:focal}.

Existen diversas formulaciones de este problema. Consideraremos tres:
la integral de Rayleigh--Sommerfeld, que proporciona una representación
espacial exacta; la integral de Fresnel, obtenida como su aproximación
paraxial; y el método del espectro angular (ASM), que expresa la solución en
el dominio de las frecuencias espaciales. Presentaremos el ASM al final porque
su estructura espectral permite identificar de manera inmediata la banda
transmitida por una pupila finita y deducir a partir de ella el límite de
resolución de Abbe.

Sea \(U\) una componente escalar monocromática conocida sobre el plano
\(z=0\). Supondremos que el semiespacio \(z>0\) es homogéneo, isótropo y libre
de fuentes, de modo que
\begin{equation}
\left(\nabla^2+k^2\right)U=0,
\qquad
k=nk_0=\frac{2\pi n}{\lambda_0}
=\frac{2\pi}{\lambda_m},
\end{equation}
donde \(\lambda_0\) y \(\lambda_m=\lambda_0/n\) son las longitudes de onda en
el vacío y en el medio, respectivamente. Bajo estas hipótesis, la propagación
entre planos puede formularse de manera exacta, dentro de la teoría escalar,
tanto en el espacio real como en el dominio de las frecuencias espaciales
\cite{goodman2005}.

\subsection{Integral de Rayleigh--Sommerfeld}

El mismo problema puede formularse en el espacio real mediante la segunda
identidad de Green. Al escoger la función de Green saliente que se anula sobre
el plano \(z=0\), la representación integral de frontera se reduce a la
primera integral de difracción de Rayleigh--Sommerfeld \cite{goodman2005}:
\begin{equation}
U(\vb r,z)
=
-\frac{1}{2\pi}
\iint_{\mathbb R^2}
U(\vb r',0)
\frac{z}{R}
\left(ik-\frac{1}{R}\right)
\frac{e^{ikR}}{R}\,
d^2\vb r',
\label{eq:rayleigh_sommerfeld}
\end{equation}
donde
\begin{equation}
R=
\sqrt{
z^2+
\left\lVert
\vb r-\vb r'
\right\rVert^2
}
\end{equation}
es la distancia entre un punto del plano inicial y el punto de observación.
Esta expresión no requiere una aproximación paraxial: es una solución exacta
del problema escalar de Helmholtz en el semiespacio siempre que el medio sea
homogéneo y libre de fuentes, el campo prescrito sobre \(z=0\) sea compatible
con dicho problema y se seleccione la solución saliente mediante la condición
de radiación de Sommerfeld. La integral de Rayleigh--Sommerfeld y el método
del espectro angular que se presentará más adelante son, por tanto, las
representaciones espacial y espectral del mismo propagador.

Cuando \(kR\gg1\) para todos los puntos de la región de integración que
contribuyen de manera apreciable, el término \(1/R\) puede despreciarse frente
a \(k\). La ecuación \eqref{eq:rayleigh_sommerfeld} adopta entonces la forma
asintótica
\begin{equation}
U(\vb r,z)
\simeq
\frac{1}{i\lambda_m}
\iint_{\mathbb R^2}
U(\vb r',0)
\frac{z}{R}
\frac{e^{ikR}}{R}\,
d^2\vb r'.
\label{eq:rayleigh_sommerfeld_asymptotic}
\end{equation}
El factor \(z/R=\cos\theta\) es el factor de oblicuidad. La condición
\(kR\gg1\), equivalente a \(R\gg\lambda_m/(2\pi)\), justifica únicamente esta
simplificación de la amplitud; no constituye por sí sola la aproximación de
Fresnel.

\subsection{Aproximación de Fresnel}

La aproximación de Fresnel se obtiene cuando las separaciones transversales
relevantes son pequeñas comparadas con la distancia de propagación. Si
\(\Delta\vb r
=\vb r-\vb r'\), entonces
\begin{equation}
R
=
z\sqrt{1+\frac{\lVert\Delta\vb r\rVert^2}{z^2}}
=
z+\frac{\lVert\Delta\vb r\rVert^2}{2z}
-\frac{\lVert\Delta\vb r\rVert^4}{8z^3}
+\cdots.
\end{equation}
En la amplitud se toman \(R\simeq z\) y \(z/R\simeq1\), mientras que en la fase
se conserva el término cuadrático. Además de
\(\lVert\Delta\vb r\rVert/z\ll1\), la omisión del término de cuarto
orden requiere que, sobre la región efectiva de integración,
\begin{equation}
\frac{k\lVert\Delta\vb r\rVert^4}{8z^3}\ll1.
\label{eq:fresnel_validity}
\end{equation}
Con estas aproximaciones, la ecuación
\eqref{eq:rayleigh_sommerfeld_asymptotic} conduce a
\begin{equation}
U(\vb r,z)
\simeq
\frac{e^{ikz}}{i\lambda_m z}
\iint_{\mathbb R^2}
U(\vb r',0)
\exp\!\left[
\frac{ik}{2z}
\left\lVert
\vb r-\vb r'
\right\rVert^2
\right]
d^2\vb r'.
\label{eq:fresnel_integral}
\end{equation}
Esta es la integral de Fresnel. Define un sistema lineal e invariante ante
traslaciones, por lo que puede escribirse como una convolución,
\begin{equation}
U(\vb r,z)
=
U(\vb r,0)
*
\left[
\frac{e^{ikz}}{i\lambda_m z}
\exp\!\left(
\frac{ik}{2z}\lVert\vb r\rVert^2
\right)
\right].
\end{equation}
Al separar los términos cuadráticos del núcleo, también se obtiene la forma
de transformada de Fourier
\begin{equation}
U(\vb r,z)
=
\frac{e^{ikz}}{i\lambda_m z}
\exp\!\left(
\frac{ik}{2z}\lVert\vb r\rVert^2
\right)
\mathscr F_\perp
\left\{
U(\vb r',0)
\exp\!\left(
\frac{ik}{2z}\lVert\vb r'\rVert^2
\right)
\right\}_{\mathbf k_\perp
=k\vb r/z}.
\label{eq:fresnel_fourier_form}
\end{equation}

La integral de Fresnel también puede recuperarse a partir del método del
espectro angular si el número de onda longitudinal se expande hasta segundo
orden en \(k_\perp/k\). Ambas vías conducen al mismo propagador paraxial.

\subsection{Método del espectro angular}

El método del espectro angular descompone el campo inicial en ondas planas. Con
\(\vb r=(x,y)\), \(\mathbf k_\perp=(k_x,k_y)\) y la convención
\begin{align}
\widetilde U_0(\mathbf k_\perp)
&=
\iint_{\mathbb R^2}
U(\vb r,0)
e^{-i\mathbf k_\perp\cdot\vb r}\,
d^2\vb r,
\\
U(\vb r,0)
&=
\frac{1}{(2\pi)^2}
\iint_{\mathbb R^2}
\widetilde U_0(\mathbf k_\perp)
e^{i\mathbf k_\perp\cdot\vb r}\,
d^2\mathbf k_\perp,
\end{align}
la ecuación de Helmholtz se reduce, para cada componente espectral, a
\begin{equation}
\frac{\partial^2\widetilde U}{\partial z^2}
+
\left(k^2-k_\perp^2\right)\widetilde U=0,
\qquad
k_\perp=\lvert\mathbf k_\perp\rvert.
\end{equation}
La solución que se propaga hacia \(z>0\) es
\begin{equation}
\widetilde U(\mathbf k_\perp,z)
=
\widetilde U_0(\mathbf k_\perp)e^{ik_z z},
\qquad
k_z=\sqrt{k^2-k_\perp^2},
\qquad
\operatorname{Im}(k_z)\geq0.
\end{equation}
La elección de rama impide el crecimiento exponencial de las componentes
evanescentes. Para \(k_\perp\leq k\), \(k_z\) es real y las componentes son
propagantes; para \(k_\perp>k\), \(k_z\) es imaginario y las componentes
decaen con la distancia.

Al aplicar la transformada inversa se obtiene el propagador del espectro
angular,
\begin{equation}
U(\vb r,z)
=
\frac{1}{(2\pi)^2}
\iint_{\mathbb R^2}
\widetilde U_0(\mathbf k_\perp)
\exp\!\left[
i\left(
\mathbf k_\perp\cdot\vb r
+k_z z
\right)
\right]
d^2\mathbf k_\perp.
\label{eq:angular_spectrum_propagation}
\end{equation}
En forma operacional,
\begin{equation}
\mathscr R_{\mathrm{ASM}}(z)
=
\mathscr F_\perp^{-1}
\left\{
e^{ik_z(\mathbf k_\perp)z}
\mathscr F_\perp\left\{\;\cdot\;\right\}
\right\}.
\label{eq:propagation_operator}
\end{equation}
Esta representación es exacta para el problema escalar planteado y resulta
especialmente conveniente para la propagación numérica mediante transformadas
rápidas de Fourier.

\subsection{Apertura numérica y límite de resolución}
\label{sec:resolution_limit}

Una pupila finita restringe el intervalo de direcciones que puede atravesar el
sistema. Si \(\alpha\) es el semiángulo máximo de aceptación, la apertura
numérica es \cite{bornwolf1999}
\begin{equation}
\mathrm{NA}=n\sin\alpha.
\label{eq:def_NA}
\end{equation}
Para una componente plana del espectro angular, la dirección de propagación
satisface
\begin{equation}
\sin\theta
=
\frac{k_\perp}{k}.
\end{equation}
Como la pupila solo admite componentes con \(\lvert\theta\rvert\leq\alpha\),
esta restricción equivale a
\begin{equation}
k_\perp\leq k\sin\alpha=k_0\mathrm{NA},
\qquad
\lvert f_\perp\rvert
\leq
\frac{\mathrm{NA}}{\lambda_0},
\qquad
f_\perp=\frac{k_\perp}{2\pi}.
\label{eq:pupil_spatial_frequency_cutoff}
\end{equation}
La pupila actúa así como un filtro que elimina las frecuencias espaciales
transversales situadas fuera de su banda de paso. Adoptando \(d\) como el
semiperiodo de una estructura periódica, de modo que su frecuencia espacial
fundamental es \(1/(2d)\), la condición de transmisión
\begin{equation}
\frac{1}{2d}
\leq
\frac{\mathrm{NA}}{\lambda_0}
\end{equation}
conduce al límite de Abbe \cite{Abbe1873},
\begin{equation}
d_{\min}
=
\frac{\lambda_0}{2\,\mathrm{NA}}.
\label{eq:abbe_resolution_limit}
\end{equation}
La definición de \(d\) es relevante: el periodo completo correspondiente es
\(2d_{\min}=\lambda_0/\mathrm{NA}\). Esta forma hace explícitas las
convenciones espectrales y evita atribuir el factor \(2\) al propagador.

El aumento de la apertura numérica extiende la banda de frecuencias
transmitidas y mejora la resolución, pero también incrementa el intervalo
angular presente en el campo. En ese régimen, las direcciones locales de
polarización y los coeficientes de Fresnel varían de manera apreciable sobre
la pupila, por lo que la descripción vectorial desarrollada a continuación
se vuelve necesaria.

\section{Transformación que sufre el campo eléctrico al cambiar de medio} 


Considérese dos medios dieléctricos $M_1$ \textit{\&} $M_2$ separados por una superficie $\Sigma$ cuya normal en cada punto se considera positiva cuando apunta hacia el medio $M_1$.

Sean $n_1$ y $n_2$ los índices de refracción de $M_1$ y $M_2$ respectivamente; el campo eléctrico $\vb E_1$ en el medio $M_1$ y el campo eléctrico $\vb E_2$ en el medio $M_2$ satisfacen
\begin{equation}
  \nabla^2 \vb E_i + k_i^2 \vb E_i = \vb 0,
  \qquad
  i=1,2,
\end{equation}
con
\[
  k_i = n_i k_0 = \frac{2\pi n_i}{\lambda_0}.
\]

Como la ecuación de Helmholtz a satisfacer es distinta en cada medio, esto genera una discontinuidad que transforma el campo que incide en la superficie $\Sigma$.


Escribiendo el campo por su espectro
\begin{equation*} 
  \vb{E}(\vb x) = \frac{1}{(2\pi)^3} \int_{\mathbb K^3} \widehat{\vb{E}} (\vb{k}) e^{i\vb k \cdot \vb x} d^3k,
\end{equation*}
que por cumplir la ecuación de Helmholtz, es equivalente a
\begin{equation*} 
  \vb{E}(\vb x) = \frac{1}{(2\pi)^3} \int_{\tilde{S}^2} \widehat{\vb{E}} (\vb{k}) e^{ik \uvec k \cdot \vb x} d\tilde{\Omega},
\end{equation*}
donde $\tilde S^2:=\{\hat{\mathbf k}\in\mathbb R^3:\|\hat{\mathbf k}\|=1\}$ es la esfera unitaria en el espacio $\mathbf{k}$, y $d\tilde\Omega$ su elemento de superficie (ángulo sólido), con $d\tilde\Omega=\sin\tilde\theta\,d\tilde\theta\,d\tilde\varphi$ en coordenadas esféricas.

Al remplazarlo en las condiciones de frontera \cite{jackson1999}, se encuentra la ley de transformación para cada componente $\vb k$ del campo.

Dicha transformación puede ser efectuada de forma matemática mediante los coeficientes de Fresnel para una onda plana y monocromática, o la componente espectral de un campo general, la deducción de dicha transformación puede encontrarse en el apéndice \ref{sec:boundary_conditions_in_dielectrics}.


El campo incidente $\widehat{\vb E}(\vb k)$ se transforma de forma diferente en cada punto de la interfaz, y de manera diferente en dos polarizaciones principales $\uvec s$ y $\uvec p$.
Para definir las polarizaciones $\uvec s$ y $\uvec p$ primero hay que definir el plano $\Pi$ de incidencia
\begin{equation}
\Pi(\uvec{k},\uvec N) = \operatorname{span}\{\uvec{k}, \uvec{N}\},
\label{eq:pi}
\end{equation}
definimos así entonces los vectores unitarios $\uvec s$ y $\uvec p$ según
\begin{equation}
  \uvec{s}(\uvec{k},\uvec N)
  = \frac{\uvec N\times\uvec{k}}
  {\left\|\uvec N\times\uvec{k}\right\|},
\label{eq:sdef}
\end{equation}

\begin{equation}
  \uvec{p}(\uvec{k},\uvec N)
  = \uvec{s}(\uvec{k},\uvec N)\times\uvec{k}
  = \frac{(\uvec N\cdot\uvec{k})\uvec{k}-\uvec N}
  {\sqrt{1-\left(\uvec N\cdot\uvec{k}\right)^2}}.
\label{eq:pdef}
\end{equation}

El vector $\uvec{s}$ es perpendicular al plano de incidencia, mientras que $\uvec{p}$ pertenece a $\Pi(\uvec{k},\uvec N)$ y es perpendicular a $\uvec{k}$. Para el campo transmitido se definen análogamente $\uvec{s}'=\uvec{s}(\uvec{k}',\uvec N)$ y $\uvec{p}'=\uvec{p}(\uvec{k}',\uvec N)$. La figura~\ref{fig:fresnel_sp_interface} ilustra esta descomposición en el caso elemental de una interfaz plana: la componente $s$ del campo sale del plano de incidencia y es común a las ondas incidente, reflejada y transmitida, mientras que la componente $p$ yace en dicho plano y gira con la dirección de propagación de cada onda.

\begin{figure}[H]
  \centering
  \thesisincludegraphics[width=0.5\textwidth]
    {figures/geometrical_setups/fresnel_ps.pdf}
    {figures/fast/fresnel_ps.png}
  \caption{Descomposición del campo eléctrico en las polarizaciones $s$ y $p$ en una interfaz plana entre dos medios de índices $n$ y $n'$. La componente $\vb E_s$ es perpendicular al plano de incidencia (sale del dibujo), y la componente $\vb E_p$ yace en él; cada onda (incidente, reflejada y transmitida) lleva su propio par $(\vb E_s, \vb E_p)$.}
  \label{fig:fresnel_sp_interface}
\end{figure}

\begin{figure}[H]
  \centering
  \thesisincludegraphics[width=0.55\textwidth]
    {figures/geometrical_setups/base_fresnel.pdf}
    {figures/fast/base_fresnel.jpg}
  \caption{Base de Fresnel en un punto $Q$ de la interfaz refractante. El plano de incidencia $\Pi = \operatorname{span}\{\hat{u}, \hat{N}\}$ contiene la normal $\hat{N}$ y la dirección de propagación incidente $\hat{u}$; el vector $\hat{p}$ (y su homólogo transmitido $\hat{p}'$) yace en $\Pi$ perpendicular a $\hat{u}$ (resp.\ $\hat{u}'$), mientras que $\hat{s}$ es perpendicular a $\Pi$.}
  \label{fig:fresnel_basis}
\end{figure}

Los coeficientes $t_s$, $t_p$ se conocen como coeficientes de Fresnel para la transmisión \cite{bornwolf1999} y se definen como
\begin{equation}
  t_s =
  \frac{2 n_1 \cos\theta_i}
  {n_1\cos\theta_i+n_2\cos\theta_t},
  \label{eq:ts}
\end{equation}

\begin{equation}
  t_p =
  \frac{2 n_1 \cos\theta_i}
  {n_2\cos\theta_i+n_1\cos\theta_t}.
  \label{eq:tp}
\end{equation}

Para escribirlos directamente en términos de la dirección del vector de onda incidente $\uvec{k}$, usamos
\begin{equation}
  \cos\theta_i = -\,\uvec N\cdot\uvec{k}.
  \label{eq:cos_i_kN}
\end{equation}
Por la ley de Snell, $n_1\sin\theta_i=n_2\sin\theta_t$, se tiene
\begin{equation}
  \cos\theta_t
  = \sqrt{1-\left(\frac{n_1}{n_2}\right)^2
  \left[1-\left(\uvec N\cdot\uvec{k}\right)^2\right]},
  \label{eq:cos_t_kN}
\end{equation}
donde se toma la rama positiva para la onda transmitida propagándose hacia el segundo medio. Sustituyendo \eqref{eq:cos_i_kN} y \eqref{eq:cos_t_kN} en \eqref{eq:ts} y \eqref{eq:tp}, los coeficientes de Fresnel de transmisión quedan como
\begin{equation}
t_s(\uvec{k},\uvec N) =
\frac{-2 n_1\,\uvec N\cdot\uvec{k}}
{-n_1\,\uvec N\cdot\uvec{k}
+ n_2\sqrt{1-\left(\frac{n_1}{n_2}\right)^2
\left[1-\left(\uvec N\cdot\uvec{k}\right)^2\right]}},
\label{eq:ts_kN}
\end{equation}

\begin{equation}
t_p(\uvec{k},\uvec N) =
\frac{-2 n_1\,\uvec N\cdot\uvec{k}}
{-n_2\,\uvec N\cdot\uvec{k}
+ n_1\sqrt{1-\left(\frac{n_1}{n_2}\right)^2
\left[1-\left(\uvec N\cdot\uvec{k}\right)^2\right]}}.
\label{eq:tp_kN}
\end{equation}

Con esto, conociendo el campo incidente sobre la interfaz $\Sigma$ que separa dos medios, podemos determinar el campo transmitido punto a punto sobre la misma interfaz. Si $\vb x_\Sigma\in\Sigma$, entonces la forma operacional de esta transformación es
\begin{equation}
  \vb{E}'_{\Sigma}(\vb x_\Sigma)
  =
  \mathscr{T}\!\left[
  \uvec{k},\uvec N(\vb x_\Sigma)
  \right]
  \left\{\vb{E}_{\Sigma}(\vb x_\Sigma)\right\}.
\end{equation}

En notación bra-ket, el operador local de transmisión se escribe como
\begin{equation}
  \mathscr{T}\!\left[
  \uvec{k},\uvec N(\vb x_\Sigma)
  \right]
  =
  t_p\!\left[
  \uvec{k},\uvec N(\vb x_\Sigma)
  \right]
  \ket{\uvec{p}'}
  \bra{\uvec{p}}
  +
  t_s\!\left[
  \uvec{k},\uvec N(\vb x_\Sigma)
  \right]
  \ket{\uvec{s}'}
  \bra{\uvec{s}}.
\end{equation}

De forma explícita, el campo transmitido queda dado por
\begin{equation}
  \begin{aligned}
  \vb{E}'_{\Sigma}(\vb x_\Sigma)
  &=
  t_p\!\left[
  \uvec{k},\uvec N(\vb x_\Sigma)
  \right]
  \left[
  \vb{E}_{\Sigma}(\vb x_\Sigma)\cdot
  \uvec{p}\!\left[
  \uvec{k},\uvec N(\vb x_\Sigma)
  \right]
  \right]
  \uvec{p}'\!\left[
  \uvec{k}',\uvec N(\vb x_\Sigma)
  \right]
  \\[4pt]
  &\quad+
  t_s\!\left[
  \uvec{k},\uvec N(\vb x_\Sigma)
  \right]
  \left[
  \vb{E}_{\Sigma}(\vb x_\Sigma)\cdot
  \uvec{s}\!\left[
  \uvec{k},\uvec N(\vb x_\Sigma)
  \right]
  \right]
  \uvec{s}'\!\left[
  \uvec{k}',\uvec N(\vb x_\Sigma)
  \right].
  \end{aligned}
  \label{eq:Transmited_E_field_in_Sigma}
\end{equation}

\textit{(véase \ref{sec:appendix_maxwell_frontera_fresnel} para una deducción detallada de este resultado)}.

En el problema \textit{estigmático simple} la dirección de propagación en función de cualquier punto $Q \ \in \Sigma$ es siempre conocida, pues por la construcción del problema estigmático
$$
\uvec k := \uvec u = \text{dirección}(A, Q),
$$

$$
\uvec k':= \uvec u' = \text{dirección}(Q, A'),
$$
donde dirección($AQ$) es el vector unitario que va desde $A$ hasta $Q$.
Esto porque todos los rayos que emergen de $A$ lo hacen radialmente, e igualmente radialmente convergen en $A'$.

Tenemos pues, todo lo necesario para pasar al estudio del campo eléctrico en el plano focal en un sistema simple refractivo.



% chapter  (end)
